% ─────────────────────────────────────────────────────────────────────────────
\section{Requirements, Constraints, and Final Scope}
\label{sec:reqs}
% ─────────────────────────────────────────────────────────────────────────────

The final system was shaped by a combination of functional expectations, operational
constraints, and practical scope limits, including object detection, deep semantic reasoning, offline execution,
structured output generation, and operator-facing presentation.

\subsection{Functional Requirements}

\begin{itemize}[nosep]
  \item \textbf{FR-1:} The system shall ingest thermal imagery from UAV or UAV-like local
    video sources.
  \item \textbf{FR-2:} The system shall detect multiple objects of interest in a frame,
    including at least people, vehicles, animals, weapons, or ammunition-related classes
    where supported by training data.
  \item \textbf{FR-3:} The system shall crop regions of interest from detector outputs and
    pass them to a second-stage reasoning component.
  \item \textbf{FR-4:} The system shall generate semantic analysis richer than plain labels,
    such as likely object type, possible status, or qualitative threat relevance.
  \item \textbf{FR-5:} The system shall produce structured machine-readable output in JSON
    form for logging and interface consumption.
  \item \textbf{FR-6:} The system shall operate without dependence on external web APIs or
    cloud inference during runtime.
  \item \textbf{FR-7:} The system shall support local presentation of alerts, detections, or
    analytical summaries for the operator.
\end{itemize}

\subsection{Non-Functional Requirements}

\begin{itemize}[nosep]
  \item \textbf{Reliability:} The pipeline should behave consistently on representative
    thermal scenes and should degrade gracefully when no object is found.
  \item \textbf{Security:} Data acquisition, processing, and storage must remain inside a
    local or air-gapped environment with zero intentional external transfer.
  \item \textbf{Performance:} The design should remain compatible with near-real-time
    operation and edge deployment goals, even if final benchmarking depends on hardware
    availability.
  \item \textbf{Usability:} The operator output should be more actionable than raw
    detections and should reduce interpretation burden rather than increase it.
  \item \textbf{Maintainability:} The architecture should remain modular so that the
    detector, reasoner, prompts, or logging layer can be updated independently.
\end{itemize}

\subsection{Project Constraints}

Several constraints influenced the final design choices:

\begin{itemize}[nosep]
  \item \textbf{Closed-network constraint:} the system could not rely on remote inference or
    hosted model APIs.
  \item \textbf{Thermal-domain constraint:} training data and model behavior are limited by
    the intrinsic ambiguity of thermal imagery.
  \item \textbf{Hardware constraint:} the target vision involved deployment on resource-
    constrained edge hardware such as the NVIDIA Jetson family.
  \item \textbf{Open-source constraint:} the project aimed to use publicly available models,
    frameworks, and datasets wherever possible.
  \item \textbf{Capstone scope constraint:} the project needed to stay feasible within a senior
    project schedule rather than expand into full-scale field qualification.
\end{itemize}

\subsection{Final Scope Boundaries}

The final scope emphasized the following boundaries are important:

\begin{itemize}[nosep]
  \item The system aims for \emph{contextual understanding}, but not guaranteed forensic-
    grade identification of exact commercial vehicle make and model.
  \item The system is a \emph{decision-support} tool and does not automate engagement or
    tactical response.
\end{itemize}

These scope boundaries are consistent with both the project's educational setting and the
real technical difficulty of thermal semantic reasoning.
